% Part: first-order-logic
% Chapter: completeness
% Section: henkin-expansion

\documentclass[../../../include/open-logic-section]{subfiles}

\begin{document}

\olfileid{fol}{com}{hen}

\olsection{Henkin 展开}

\begin{explain}
证明完备性定理的挑战之一在于，我们从完备一致集~$\Gamma$ 构造出的模型，必须使~$\Gamma$ 中的所有量词公式为真。为了保证这一点，我们使用一个源于 Leon Henkin（莱昂·亨金）的技巧。本质上，这个技巧在于通过添加可数无穷多个新的常项符号来扩充语言，并且为每一个含有一个自由变元 $x$ 的公式 $!A(x)$，添加一个形如
\iftag{prvEx}
      {$\lexists[x][!A(x)] \lif !A(c)$}
      {$\lnot\lforall[x][!A(x)] \lif \lnot !A(c)$}
的公式，其中 $c$ 是某个新的常项符号。当我们构造满足~$\Gamma$ 的结构时，这将保证每个
\iftag{prvEx}
{真的存在语句在新常项中都有一个见证}
{假的全称语句在新常项中都有一个反例}。
\end{explain}

\begin{prop}
\ollabel{prop:lang-exp}
设 $\Gamma$ 在语言 $\Lang L$ 中是一致的，且 $\Lang L'$ 是通过给 $\Lang L$ 添加一个可数无穷的新常项符号集 $\Obj d_0$,
$\Obj d_1$, \dots 而得到的，则 $\Gamma$ 在~$\Lang L'$ 中也是一致的。
\end{prop}

\begin{defn}[饱和集]
语言 $\Lang{L}$ 的一个公式集 $\Gamma$ 是\emph{饱和的}，当且仅当对于每一个含有一个自由变元~$x$ 的公式~$!A(x) \in \Frm[L]$，都存在一个常项符号~$c \in \Lang{L}$，使得
\iftag{prvEx}
      {$\lexists[x][!A(x)] \lif !A(c) \in \Gamma$}
      {$\lnot\lforall[x][!A(x)] \lif \lnot !A(c) \in \Gamma$}。
\end{defn}

以下定义将用于下一定理的证明。

\begin{defn}
\ollabel{defn:henkin-exp}
设 $\Lang L'$ 如 \olref{prop:lang-exp} 所述。固定~$\Lang L'$ 中所有恰有一个变元（$x_i$）自由出现的公式~$!A_i(x_i)$ 的一个枚举 $!A_0(x_0)$, $!A_1(x_1)$, \dots。我们归纳地定义语句~$!D_n$。

令 $c_0$ 为添加到~$\Lang{L}$ 的新常项符号 $\Obj d_i$ 中，第一个不在~$!A_0(x_0)$ 中出现的常项符号。假设 $!D_0$, \dots,~$!D_{n-1}$ 已经定义，令 $c_n$ 为新常项符号~$\Obj d_i$ 中第一个既不出现在 $!D_0$, \dots,~$!D_{n-1}$ 中，也不出现在~$!A_n(x_n)$ 中的常项符号。

现令 $!D_{n}$ 为公式 \iftag{prvEx}
{$\lexists[x_{n}][!A_{n}(x_{n})] \lif
  !A_{n}(c_{n})$}{$\lnot\lforall[x_{n}][!A_{n}(x_{n})] \lif \lnot
  !A_{n}(c_{n})$}。
\end{defn}

\begin{lem}
\ollabel{lem:henkin}
每个一致集~$\Gamma$ 都可以扩充为一个饱和的一致集~$\Gamma'$。
\end{lem}

\begin{proof}
给定语言~$\Lang{L}$ 中的一个一致语句集~$\Gamma$，通过添加一个可数无穷的新常项符号集来扩充语言，得到~$\Lang{L'}$。根据 \olref{prop:lang-exp}，$\Gamma$ 在更丰富的语言中仍是一致的。进一步，令 $!D_i$ 如 \olref{defn:henkin-exp} 所定义。令
\begin{align*}
\Gamma_0 & = \Gamma \\
\Gamma_{n+1} & = \Gamma_n \cup \{!D_n \}
\end{align*}
即 $\Gamma_{n+1} = \Gamma \cup \{ !D_0, \dots, !D_n \}$，并令
$\Gamma' = \bigcup_{n} \Gamma_n$。显然 $\Gamma'$ 是饱和的。

如果 $\Gamma'$ 不一致，那么对于某个 $n$，$\Gamma_n$ 将会不一致（习题：解释原因）。因此，要证明 $\Gamma'$ 是一致的，只需对~$n$ 进行归纳，证明每个集合~$\Gamma_n$ 是一致的即可。

归纳基础即 $\Gamma_0 = \Gamma$ 是一致的，这正是定理的假设。对于归纳步骤，假设 $\Gamma_{n}$ 是一致的但 $\Gamma_{n+1} =
\Gamma_n \cup \{!D_n\}$ 不一致。回顾 $!D_n$~是
\iftag{prvEx}
{$\lexists[x_{n}][!A_{n}(x_n)] \lif !A_{n}(c_{n})$}
{$\lnot\lforall[x_{n}][!A_{n}(x_n)] \lif \lnot !A_{n}(c_{n})$}，
其中 $!A_n(x_n)$ 是 $\Lang{L'}$ 中仅变元~$x_n$ 是自由的公式。根据我们选择~$c_n$ 的方式（见 \olref{defn:henkin-exp}），$c_n$ 不出现在~$!A_n(x_n)$ 中，也不出现在~$\Gamma_n$ 中。

如果 $\Gamma_n \cup \{!D_n\}$ 不一致，则 $\Gamma_n
\Proves \lnot !D_n$，从而以下两者同时成立：
\iftag{prvEx}{
\[
\Gamma_n \Proves \lexists[x_n][!A_n(x_n)]
\qquad
\Gamma_n \Proves \lnot !A_n(c_n)
\]}{
\[
\Gamma_n \Proves \lnot\lforall[x_n][!A_n(x_n)]
\qquad
\Gamma_n \Proves !A_n(c_n)
\]}
由于 $c_n$ 不出现在
$\Gamma_n$ 或~$!A_n(x_n)$ 中，
\tagrefs{prfAX/{fol:axd:qpr:thm:strong-generalization},prfSC/{fol:seq:qpr:thm:strong-generalization},prfND/{fol:ntd:qpr:thm:strong-generalization},prfTab/{fol:tab:qpr:thm:strong-generalization}} 适用。
由 \iftag{prvEx}
{$\Gamma_n \Proves \lnot !A_n(c_n)$}
{$\Gamma_n \Proves !A_n(c_n)$}，
可得
\iftag{prvEx}
{$\Gamma_n \Proves \lforall[x_n][\lnot !A_n(x_n)]$}
{$\Gamma_n \Proves \lforall[x_n][!A_n(x_n)]$}。
于是我们有
\iftag{prvEx}
{$\Gamma_n \Proves \lexists[x_n][!A_n(x_n)]$ 与
$\Gamma_n \Proves \lforall[x_n][\lnot !A_n(x_n)]$}
{$\Gamma_n \Proves \lnot\lforall[x_n][!A_n(x_n)]$ 与
$\Gamma_n \Proves \lforall[x_n][!A_n(x_n)]$}，
因此 $\Gamma_n$ 本身不一致。
\iftag{prvEx}
{(注意
\iftag{prvAll}
{$\lforall[x_n][\lnot !A_n(x_n)] \Proves
  \lnot\lexists[x_n][!A_n(x_n)]$}
{$\lforall[x_n][\lnot !A_n]$ 定义为
  $\lnot\lexists[x_n][\lnot\lnot !A_n(x_n)]$ 与
  $\lnot\lexists[x_n][\lnot\lnot !A_n(x_n)] \Proves
  \lnot\lexists[x_n][!A_n(x_n)]$}。)}{}
矛盾：$\Gamma_n$ 本应是一致的。故
$\Gamma_n \cup \{ !D_n\}$ 是一致的。
\end{proof}

\begin{explain}
我们现在要说明，\emph{完备}、一致且饱和的集合具有这样的性质：\iftag{prvAll}{它包含一个全称量化语句当且仅当它包含其所有实例}{}\iftag{defAll,defEx}{}{并且}\iftag{prvAll}{它包含一个存在量化语句当且仅当它包含至少一个实例}{}。我们将利用这一点来证明，从一个完备、一致、饱和的集合生成的结构，会使其中所有的量词语句为真。
\end{explain}

\begin{prop}\ollabel{prop:saturated-instances}
假设 $\Gamma$ 是完备、一致且饱和的。
\begin{tagenumerate}{prvEx,prvAll}
\tagitem{prvEx}{$\lexists[x][!A(x)] \in \Gamma$ 当且仅当存在至少一个闭项~$t$ 使得 $!A(t) \in \Gamma$。}{}
\tagitem{prvAll}{$\lforall[x][!A(x)] \in \Gamma$ 当且仅当对所有闭项~$t$ 都有 $!A(t) \in \Gamma$。}{}
\end{tagenumerate}
\end{prop}

\begin{proof}
  \begin{tagenumerate}{prvEx,prvAll}
  \tagitem{prvEx}{%
    \iftag{probEx}{练习。}{首先假设 $\lexists[x][!A(x)]
      \in \Gamma$。因为 $\Gamma$ 是饱和的，
      所以对某个常项符号~$c$ 有 $(\lexists[x][!A(x)] \lif !A(c)) \in \Gamma$。根据
      \tagrefs{prfAX/{fol:axd:ppr:prop:provability-lif},prfSC/{fol:seq:ppr:prop:provability-lif},prfND/{fol:ntd:ppr:prop:provability-lif},prfTab/{fol:tab:ppr:prop:provability-lif}} 的第~(1) 条，
      以及 \olref[ccs]{prop:ccs}\olref[ccs]{prop:ccs-prov-in}，可得 $!A(c)
      \in \Gamma$。

    反之，不需要饱和性：假设
    $!A(t) \in \Gamma$。则由
      \tagrefs{prfAX/{fol:axd:qpr:prop:provability-quantifiers},prfSC/{fol:seq:qpr:prop:provability-quantifiers},prfND/{fol:ntd:qpr:prop:provability-quantifiers},prfTab/{fol:tab:qpr:prop:provability-quantifiers}} 的第~(1) 条，可得 $\Gamma \Proves \lexists[x][!A(x)]$。再根据
    \olref[ccs]{prop:ccs}\olref[ccs]{prop:ccs-prov-in}，
    得到 $\lexists[x][!A(x)] \in \Gamma$。}}{}
  \tagitem{prvAll}{%
    \iftag{probAll}{练习。}{假设对所有闭项~$t$ 都有 $!A(t) \in \Gamma$。用反证法，假定
      $\lforall[x][!A(x)] \notin \Gamma$。由于 $\Gamma$ 是完备的，
      因此 $\lnot\lforall[x][!A(x)] \in \Gamma$。根据饱和性，
      对某个常项符号~$c$ 有 \iftag{prvEx}{$(\lexists[x][\lnot !A(x)] \lif \lnot !A(c)) \in
        \Gamma$}{$(\lnot\lforall[x][!A(x)] \lif \lnot !A(c)) \in
        \Gamma$}。由假设，因为 $c$
      是闭项，所以 $!A(c) \in \Gamma$。但这将导致
      $\Gamma$ 不一致。（练习：给出推导以证明
      \[
      \lnot \lforall[x][!A(x)],
      \iftag{prvEx}{\lexists[x][\lnot !A(x)] \lif \lnot
        !A(c)}{\lnot\lforall[x][!A(x)] \lif \lnot
        !A(c)}, !A(c)
      \]
      是不一致的。）

      反之，我们不需要饱和性：假设
      $\lforall[x][!A(x)] \in \Gamma$。那么，由
      \tagrefs{prfAX/{fol:axd:qpr:prop:provability-quantifiers},prfSC/{fol:seq:qpr:prop:provability-quantifiers},prfND/{fol:ntd:qpr:prop:provability-quantifiers},prfTab/{fol:tab:qpr:prop:provability-quantifiers}}
      的第~(2) 条，有 $\Gamma \Proves !A(t)$。再根据
      \olref[ccs]{prop:ccs}，得到 $!A(t) \in \Gamma$。}}{}
  \end{tagenumerate}
\end{proof}

\end{document}